\subsection{Multiplikation mit einem Skalar}
\label{subsec:vector_algebra_scalar_multiplication}

Die erste Mulitplikation, die wir in der Vektorrechnung kennenlernen,
ist die Mulitplikation eines Skalars mit einem Vektor.

\begin{DefinitionBox}{Multiplikation mit einem Skalar}{vector_algebra_scalar_multiplication}
  Die \emph{Multiplikation mit einem Skalar} \EnglishTerm{scalar multiplication}
  erfolgt für $\lambda\in\R$ und $\vec a\in\R^n$ komponentenweise:
  \[
    \lambda\vec a
    =\begin{ClozeContinuation}[15]{5}
      \lambda\begin{pmatrix}a_1\\a_2\\\vdots\\a_n\end{pmatrix}
      =\begin{pmatrix}\lambda a_1\\\lambda a_2\\\vdots\\\lambda a_n\end{pmatrix}
    \end{ClozeContinuation}.
  \]
  Das Ergebnis $\lambda\vec a\in\R^n$ ist wieder ein Vektor.
\end{DefinitionBox}

\begin{ExampleBox}{Einen Vektor mit einem Skalar multiplizieren}{vector_algebra_scale_example}
  \begin{minipage}[c]{0.49\linewidth}
    Gegeben sei der Vektor
    \[
      {\color{black}
        \vec v=\begin{pmatrix}2\\1\end{pmatrix}}.
    \]
    Die Multiplikation mit dem Skalar $\lambda$ für
    \begin{itemize}[leftmargin=1.4em,itemsep=0.8em]
      \item {\color{olive!80!black}
        $\lambda=2$ entspricht einer \Cloze{2.8cm}{Streckung}:}
        \[
          {\color{olive!80!black}2\vec v=}
          \begin{ClozeContinuation}[11]{2}
            {\color{olive!80!black}
              2\begin{pmatrix}2\\1\end{pmatrix}
              =\begin{pmatrix}4\\2\end{pmatrix}}
          \end{ClozeContinuation}.
        \]
      \item {\color{THMBlue}
        $\lambda=-1$ liefert den \Cloze{3.1cm}{Gegenvektor}:}
        \[
          {\color{THMBlue}(-1)\vec v=}
          \begin{ClozeContinuation}[12]{2}
            {\color{THMBlue}
              (-1)\begin{pmatrix}2\\1\end{pmatrix}
              =\begin{pmatrix}-2\\-1\end{pmatrix}}
          \end{ClozeContinuation}.
        \]
    \end{itemize}
  \end{minipage}\hfill
  \begin{minipage}[c]{0.48\linewidth}
    \centering
    \CachedTikzFigure{content/resources/chapters/vector_algebra/tikz/scalar_multiplication_example_2d.tex}
  \end{minipage}
\end{ExampleBox}

Die geometrische Wirkung der Multiplikation mit einem Skalar lässt sich
folgendermaßen zusammenfassen:

\begin{KeyPointBox}{Geometrische Wirkung der Multiplikation mit einem Skalar}{vector_algebra_scalar_effect}
  \begin{center}
    \CachedTikzFigure{content/resources/chapters/vector_algebra/tikz/scalar_multiplication_cases.tex}
  \end{center}
  \[
    \text{Für }
    \begin{gathered}
      \lambda>0\\[0.5em]
      \lambda<0
    \end{gathered}
    \text{ ist der resultierende Vektor }\lambda\vec a\text{ }
    \begin{gathered}
      \text{parallel}\\[0.5em]
      \text{antiparallel}
    \end{gathered}
    \text{ zu }\vec a,
    \qquad\text{d.\,h.}\qquad
    \begin{gathered}
      \lambda\vec a\upuparrows\vec a,\\[0.5em]
      \lambda\vec a\updownarrows\vec a.
    \end{gathered}
  \]
\end{KeyPointBox}

Da in jeder Komponente ausschließlich Zahlen multipliziert werden, übertragen
sich die entsprechenden Rechengesetze der Zahlenmultiplikation auf die
Multiplikation eines Vektors mit Skalaren.

\begin{TheoremBox}{Rechenregeln der Multiplikation mit einem Skalar}{vector_algebra_scalar_laws}
  Für $\alpha,\beta\in\R$ und $\vec a,\vec b\in\R^n$ gilt
  \[
    \begin{array}{@{}c@{\qquad}l@{}}
      \text{--}
        & \text{Kommutativgesetz}\\[1.2em]
      \alpha(\beta\vec a)=(\alpha\beta)\vec a
        & \text{Assoziativgesetz}\\[1.2em]
      \left.
        \begin{aligned}
          (\alpha+\beta)\vec a&=\alpha\vec a+\beta\vec a,\\[0.4em]
          \alpha(\vec a+\vec b)&=\alpha\vec a+\alpha\vec b
        \end{aligned}
        \hspace{0.8em}
      \right\}
        & \text{Distributivgesetze}
    \end{array}
  \]
\end{TheoremBox}

Daher kann bei der Multiplikation eines Vektors mit Skalaren so gerechnet
werden, wie man es von Zahlen gewohnt ist.
